% =====================================================================
%  28.4  Counting distinct shapes: interned ids as isomorphism classes.
%  Exports: sec:c22-distinct  def:c22-shaindex  thm:c22-classes
%  lst:c22-classes  fig:c22-classes
%  WORKED section only -- no solvedproblem (the chapter's two labs are 28.2
%  CF 601D and 28.3 CF 1800G).  USES 28.1's labeller as a black box; the
%  small-to-large merge it needs for the extension is Ch.17's (\chref{c28})
%  and is cited, not re-taught.  Root-dependence of the class multiset is
%  the new observation here (R has 4 rooted shape classes, its twin R'
%  rooted at 1 has 6 -- same unrooted tree).
% =====================================================================
\section{Counting Distinct Shapes: Interned Ids Are Isomorphism Classes}\label{sec:c22-distinct}

\begin{previouslybox}[Equality was the job; now the same ids become a census]
Sections 28.1--28.3 answer \emph{yes/no} questions: is \(\subtree(u)\) the shape of
\(\subtree(w)\)?  Are these two trees the same tree?  Can the children pair into
mirrors?  In every case the interned id \(\mathrm{cl}(u)\) was a \emph{key} thrown
at a compare.  But a key that names an object can also \emph{count} objects: the
map \(\mathit{tid}\) is, by construction, the list of distinct shapes that occur,
and the array of \(\mathrm{cl}\)-values is a partition of the vertices into
\term{isomorphism classes}.  Reading that partition is this whole section -- and
the reading produces one genuinely surprising fact: the census depends on the
root, even when the tree does not.
\end{previouslybox}

\begin{topicintro}
\textbf{Prerequisites.} \(\secref{c22-ahu}\) (the labeller and its shared table)
is assumed running; Chapter~17 (\chref{c28}) owns the small-to-large set merge
that the last hardening item uses.

\textbf{After this section you can}
\begin{itemize}[nosep]
  \item read off "how many different subtree shapes does this tree have" as \(\mathit{tid}.\mathrm{size()}\) after one labelling; \hfill\lv{4}
  \item answer, for every vertex, "how many vertices have a subtree isomorphic to mine" in \(O(1)\) from a frequency array; \hfill\lv{4}
  \item explain why that census is a property of the \emph{rooted} tree and changes when the root moves; \hfill\lv{4}
  \item extend the census downward -- "how many distinct shapes occur inside \(\subtree(v)\)" -- with Ch.~17 merging. \hfill\lv{5}
\end{itemize}
\end{topicintro}

% ---------------------------------------------------------------
\subsection{The partition, and what it counts}
% ---------------------------------------------------------------

\begin{definition}{Shape index and multiplicity}{c22-shaindex}
After one labelling of a rooted tree \(T\), the fibre
\([i]=\{u:\mathrm{cl}(u)=i\}\) is an \term{isomorphism class} of rooted subtrees:
exactly the vertices whose subtrees share one shape.  Write
\(m(i)=\bigl| [i]\bigr|\) for its \term{multiplicity}.  The set of occurring ids
is \(\mathit{tid}\)'s key set, so
\[
   \#\{\text{distinct subtree shapes of }T\}\;=\;|\mathit{tid}|\;=\;\#\{i: m(i)>0\}.
\]
\end{definition}

\begin{theorem}{One labelling, every census query in \(O(1)\)}{c22-classes}
Let \(T\) be rooted and labelled once with the interned-tuple fold
(\(\thmref{c22-ahu}\)).  Then, at total cost \(\bigO{n\log n}\) time and
\(\bigO{n}\) space: (i) the number of distinct shapes among the \(n\) rooted
subtrees is \(|\mathit{tid}|\); (ii) for any \(v\), the number of vertices whose
subtree is isomorphic to \(\subtree(v)\) \emph{including} \(v\) itself is
\(m(\mathrm{cl}(v))\), and excluding it, \(m(\mathrm{cl}(v))-1\); (iii) two
subtrees are isomorphic iff their ids match, so a "largest class" query is a
max over \(m\).
\end{theorem}

\begin{lstlisting}[caption={Census of subtree shapes: after the 28.1 labelling, one counting pass gives the number of distinct shapes and every class size; queries are O(1)}, label={lst:c22-classes}]
#include <bits/stdc++.h>
using namespace std;
const int MAXN=200005;
int n, nextId=1, cl[MAXN], freq[2*MAXN];
vector<int> adj[MAXN], kids[MAXN];
map<vector<int>,int> tid;                 // the shared table of 28.1
void build(int u, int p){
    for(int v: adj[u]) if(v!=p){ kids[u].push_back(v); build(v,u); }
}
int label(int u){
    vector<int> t;
    for(int v: kids[u]) t.push_back(label(v));
    sort(t.begin(), t.end());
    map<vector<int>,int>::iterator it=tid.find(t);
    if(it!=tid.end()) return cl[u]=it->second;
    int id=nextId++; tid[t]=id; return cl[u]=id;
}
int main(){
    scanf("%d",&n);
    for(int i=1;i<n;i++){ int u,v; scanf("%d%d",&u,&v); adj[u].push_back(v); adj[v].push_back(u); }
    build(1,0); label(1);
    for(int u=1;u<=n;u++) freq[cl[u]]++;                 // multiplicities
    printf("%d\n",(int)tid.size());                      // distinct shapes
    for(int u=1;u<=n;u++)                                // u, its shape id, class size
        printf("%d %d %d\n", u, cl[u], freq[cl[u]]);
    return 0;
}
\end{lstlisting}

\subsection{The census on the running tree -- and how badly the root distorts it}

\begin{figure}[H]
\centering
\begin{tikzpicture}[diagram, scale=0.95, transform shape,
   leaf/.style={circle, draw=primaryblue!70, fill=primaryblue!12, line width=0.8pt,
                minimum size=6.4mm, font=\scriptsize\bfseries, inner sep=1pt},
   star/.style={circle, draw=successgreen!70, fill=successgreen!14, line width=1pt,
                minimum size=6.4mm, font=\scriptsize\bfseries, inner sep=1pt},
   chain/.style={circle, draw=difforange!80, fill=difforange!16, line width=1pt,
                minimum size=6.4mm, font=\scriptsize\bfseries, inner sep=1pt},
   thero/.style={circle, draw=softpurple!80, fill=softpurple!14, line width=1.2pt,
                minimum size=6.4mm, font=\scriptsize\bfseries, inner sep=1pt}]
  \node[thero] (n1) at (0,0) {1};
  \node[star]  (n2) at (-3.2,-1.5) {2};
  \node[star]  (n3) at (0,-1.5) {3};
  \node[chain] (n4) at (3.2,-1.5) {4};
  \node[leaf]  (n5) at (-4.0,-3.0) {5};
  \node[leaf]  (n6) at (-2.4,-3.0) {6};
  \node[leaf]  (n7) at (-0.8,-3.0) {7};
  \node[leaf]  (n8) at (0.8,-3.0) {8};
  \node[leaf]  (n9) at (3.2,-3.0) {9};
  \draw (n1)--(n2) (n1)--(n3) (n1)--(n4) (n2)--(n5) (n2)--(n6) (n3)--(n7) (n3)--(n8) (n4)--(n9);
  \node[annot, font=\scriptsize, anchor=west, text width=4.1cm] at (4.1,0.1)
    {\(4\) classes on \(9\) vertices: blue \(\times5\), green \(\times2\), orange \(\times1\), purple \(\times1\)};
\end{tikzpicture}
\caption{The shape census of \(R\) rooted at \(1\), coloured by isomorphism class.
Class sizes: \(\mathrm{cl}=1\) (leaf) has multiplicity \(5\); \(\mathrm{cl}=2\)
(a node above two leaves) \(2\); \(\mathrm{cl}=3\) (a node above one leaf) \(1\);
\(\mathrm{cl}=4\) (the whole tree) \(1\).  So \(R\) has \(|\mathit{tid}|=4\)
distinct rooted subtree shapes, the largest class is the leaves, and every
vertex answers "how many vertices share my shape?" with one array read.}
\label{fig:c22-classes}
\end{figure}

\begin{tracebox}[The same tree, two roots: \(4\) classes against \(6\)]
\(R\) rooted at \(1\) gives the table above.  Now label the \emph{twin} \(R'\) --
the very same unrooted tree, rooted at its vertex \(1\), which is \(R\)'s vertex
\(5\) (\(\sigma^{-1}(1)=5\)) -- continuing on the same table
(\(\secref{c22-ahu}\)'s figure walks this out): leaves \(6,7,8,9\) recall
\(\mathrm{cl}=1\); vertex \(3\) recalls \(2\); vertex \(4\) recalls \(3\); vertex
\(2\) sees children \((2,3)\), a \textbf{fresh} tuple \(\to4^{\text{new}}=\mathrm{cl}\,5\);
vertex \(5\) sees \((1,5)\to\mathrm{cl}\,6\); the root \(1\) sees \((6)\to\mathrm{cl}\,7\).
Class sizes now read \(4,1,1,1,1,1\) over \(\mathbf{6}\) distinct shapes -- twice
the count for the same tree.  \textbf{Nothing went wrong}: the census counts
shapes of \emph{rooted} subtrees, and changing the root changes which subsets are
subtrees and how they hang.  A problem that asks for "distinct shapes of
subtrees" without naming a root is asking about an unrooted object, and you must
either canonicalize per component (\(\secref{c22-centre}\)) or state the root
you assumed. \(\Box\)
\end{tracebox}

\begin{insightbox}[Which counting question is which tool]
\begin{center}\renewcommand{\arraystretch}{1.25}\setlength{\tabcolsep}{3pt}
\begin{tabular}{@{}p{0.33\linewidth} p{0.28\linewidth} p{0.30\linewidth}@{}}
\toprule
\textbf{Counting question} & \textbf{Key being counted} & \textbf{Structure}\\
\midrule
how many distinct subtree shapes? & \(\mathrm{cl}(u)\) & \(|\mathit{tid}|\), or \texttt{set<int>} of ids \\
how many vertices share \(v\)'s shape? & \(\mathrm{cl}(u)\) & \texttt{freq[cl[v]]} after one pass \\
how many distinct shapes \emph{inside} \(\subtree(v)\)? & sets of ids & small-to-large merge (Ch.~17) \\
how many distinct downward \emph{strings}? & \(G\)-hashes & \(\mathcal{S}(v)\), solved CF 601D (\(\secref{c22-hash}\)) \\
how many distinct \emph{unrooted} component shapes? & centre-rooted \(\mathrm{cl}\) & \(O(1)\) compares per candidate (\(\probref{c22-summit}\)) \\
\bottomrule
\end{tabular}\end{center}
Row 4 is the trap row: strings are \emph{ordered} objects, so they need the
positional hash of \(\defref{c22-strhash}\), while rows 1--3 and 5 count
\emph{unordered} shapes and need the commutative fold or interned tuples.  Same
"count distinct" wording, two different keys.
\end{insightbox}

\begin{warningbox}[Three census errors]
\begin{enumerate}[nosep,leftmargin=1.6em]
  \item \textbf{Counting vertices instead of classes.} \(\sum_u 1 = n\) is not a
        distinct-shape count; you want \(\#\{\mathrm{cl}(u)\}\), i.e.\
        \(|\mathit{tid}|\).  Symptom: the answer never drops below \(n\) on
        trees that plainly repeat shapes.
  \item \textbf{Comparing multiplicities across tables.} \(m(i)\) is meaningful
        only inside one table: after re-interning per test case, id \(3\) in case
        A and id \(3\) in case B name different shapes.  Aggregate \emph{within}
        a case, or key everything by a collision-aware hash (\(\secref{c22-hash}\)).
  \item \textbf{Assuming the census is root-free.} The trace above shows the same
        tree yielding \(4\) and \(6\) classes.  If a problem says "subtrees of an
        unrooted tree", it usually means components after deleting a vertex -- a
        \emph{different} family of objects, handled by the centre-edge view of
        \(\secref{c22-centre}\) and the Summit's reroot pass, not by one root-at-1
        labelling.
\end{enumerate}
\end{warningbox}

\begin{hardeningbox}
\begin{itemize}[nosep,leftmargin=1.4em]
  \item \textbf{Per-subtree censuses, the cheap way.} "For every \(v\), how many
        distinct shapes occur among the subtrees inside \(\subtree(v)\)?" is a
        set-union-per-node job: give each vertex the \(\mathit{set}\) of ids in its
        subtree, merge children small-to-large, and take sizes
        (\(\bigO{n\log n}\), Ch.~17 verbatim).  The same skeleton solved CF 601D
        with strings instead of ids -- reuse your code, change the key.
  \item \textbf{Hashes for counting are \(\bigO{1}\) but not exact.}
        \texttt{unordered\_set<ull>} of \(h(u)\) instead of a \texttt{set<int>} of
        \(\mathrm{cl}\) is faster and needs no table, and it never splits one
        shape into two counts (equal shape always gives equal hash); the exposure
        is the opposite: two different shapes collapse into one count, so the
        answer can only \emph{under-count}
        (\(\thmref{c22-collision}\) -- the YES-risky direction reappearing as a
        counting error).  Prefer ids when the count is the printed answer.
  \item \textbf{Extra payload means extra key components.} Coloured or weighted
        trees ("isomorphic \emph{and} equal values") are handled by interning the
        tuple of \((\text{child ids})\) \emph{together with} the vertex value --
        i.e.\ \(\tau_u\) becomes a tuple of pairs.  One change, and every census
        query above now counts value-respecting classes.
\end{itemize}
\end{hardeningbox}

\begin{recallbox}[28.4 in one breath]
\emph{A name is a class:} one AHU labelling partitions the vertices, so
"how many shapes" is \(|\mathit{tid}|\), "who shares my shape" is
\(\mathit{freq}[\mathrm{cl}(v)]\), and both cost \(O(1)\) after
\(\bigO{n\log n}\) (\(\thmref{c22-classes}\), \(\coderef{c22-classes}\)).  On \(R\):
four classes with sizes \(5,2,1,1\) (figure \(\figref{c22-classes}\)); on the twin
\(R'\) rooted at \(1\): six classes -- because the census counts \emph{rooted}
subtrees, so state or canonicalize the root before counting.  Per-subtree
censuses are Ch.~17 merging with ids as the payload, and coloured variants just
add the colour to the interned tuple.
\selftest{
  \item One labelling done. Cost of "how many vertices have my subtree shape?" for all \(n\) vertices? \emph{(O(n) counting pass into a frequency array, then O(1) per vertex.)}
  \item Same tree, different root: same number of distinct subtree shapes? \emph{(Not necessarily -- R gives 4 classes rooted at 1 and 6 for its twin rooted at its leaf. Rooted subtrees are root-dependent objects.)}
  \item Why is "distinct downward strings" not a row of the id census? \emph{(Strings are ordered objects: their key is the positional depth-weighted hash, not the commutative shape id.)}
}
\end{recallbox}
